\documentclass[12pt]{amsart}
%prepared in AMSLaTeX, under LaTeX2e
\addtolength{\oddsidemargin}{-.65in} 
\addtolength{\evensidemargin}{-.65in}
\addtolength{\topmargin}{-.55in}
\addtolength{\textwidth}{1.3in}
\addtolength{\textheight}{.9in}

\renewcommand{\baselinestretch}{1.05}

\usepackage{verbatim} % for "comment" environment

\usepackage{palatino}

\newtheorem*{thm}{Theorem}
\newtheorem*{defn}{Definition}
\newtheorem*{example}{Example}
\newtheorem*{problem}{Problem}
\newtheorem*{remark}{Remark}

\usepackage{fancyvrb,xspace,dsfont}

\usepackage[final]{graphicx}

% macros
\usepackage{amssymb}

\usepackage{hyperref}
\hypersetup{pdfauthor={Ed Bueler},
            pdfcreator={pdflatex},
            colorlinks=true,
            citecolor=blue,
            linkcolor=red,
            urlcolor=blue,
            }

\newcommand{\br}{\mathbf{r}}
\newcommand{\bv}{\mathbf{v}}
\newcommand{\bx}{\mathbf{x}}
\newcommand{\by}{\mathbf{y}}

\newcommand{\cB}{\mathcal{B}}
\newcommand{\cD}{\mathcal{D}}
\newcommand{\cF}{\mathcal{F}}
\newcommand{\cH}{\mathcal{H}}
\newcommand{\cL}{\mathcal{L}}
\newcommand{\cV}{\mathcal{V}}
\newcommand{\cW}{\mathcal{W}}

\newcommand{\CC}{\mathbb{C}}
\newcommand{\NN}{\mathbb{N}}
\newcommand{\RR}{\mathbb{R}}
\newcommand{\ZZ}{\mathbb{Z}}

\renewcommand{\Im}{\mathrm{Im}}
\renewcommand{\Re}{\mathrm{Re}}

\newcommand{\eps}{\epsilon}
\newcommand{\grad}{\nabla}
\newcommand{\lam}{\lambda}
\newcommand{\lap}{\triangle}

\newcommand{\ip}[2]{\ensuremath{\left<#1,#2\right>}}

\newcommand{\image}{\operatorname{im}}
\newcommand{\onull}{\operatorname{null}}
\newcommand{\rank}{\operatorname{rank}}
\newcommand{\range}{\operatorname{range}}
\newcommand{\trace}{\operatorname{tr}}
\newcommand{\Span}{\operatorname{span}}

\newcommand{\prob}[1]{\bigskip\noindent\textbf{#1.}\quad }

\newcommand{\pts}[1]{(\emph{#1 pts}) }
\newcommand{\epart}[1]{\medskip\noindent\textbf{(#1)}\quad }
\newcommand{\ppart}[1]{\,\textbf{(#1)}\quad }

\newcommand{\ds}{\displaystyle}

\newcommand{\nex}{\medskip\noindent}


\begin{document}
\scriptsize \noindent Math 617 Functional Analysis (Bueler) \hfill \emph{assigned 12 January 2026; version 1.1}
\normalsize\medskip

\Large\centerline{\textbf{Assignment 1}}
\large
\medskip

\centerline{\textbf{Due Wednesday 21 January}}
\medskip
\normalsize

\thispagestyle{empty}

\bigskip
\noindent This Assignment is based on lectures and on Chapter 1 of our textbook.\footnote{Saxe (2010)~\emph{Beginning Functional Analysis}, Springer.}


\bigskip
\noindent \textsc{Do the following} Exercises for Chapter 1 (pages 13--15):

\begin{itemize}
\item \textbf{Exercise 1.1.1}
\item \textbf{Exercise 1.1.3}
\item \textbf{Exercise 1.1.5}
\item \textbf{Exercise 1.1.10}
\item \textbf{Exercise 1.3.2}
\item \textbf{Exercise 1.3.3}
\end{itemize}


\medskip
\noindent \textsc{Do the following additional problems.}

\prob{P1}  Consider the real vector space $C([0,1])$ and the linear functional
	$$\ell[f] = \int_0^1 f(x) \,e^x\, dx,$$
that is, $\ell : C([0,1]) \to \RR$.

\epart{a} Show $\ell$ is linear.

\medskip
\noindent  For a sequence of functions $\{f_k(x)\}_{k=1}^\infty \subset C([0,1])$ we define the real sequence $\ell_k = \ell[f_k] \in \RR$.

\epart{b} For $k\in \NN$ compute $\ell_k$ if $f_k(x) = x^k$.

\epart{c} For $k\in \NN$ compute $\ell_k$ if $f_k(x) = \sin(k\pi x)$.

\bigskip
\noindent  \emph{Comment (revised).}  For a finite linear combination, $g(x)=\sum_{k=1}^n c_k f_k(x)$, we have $\ell[g] = \sum_{k=1}^n c_k \ell_k$ by linearity, and from part \textbf{(b)} this defines $\ell$ on all polynomials.  If $\{f_k(x)=x^k\}$ were a ``basis'' then the values $\{\ell_k = \ell[f_k]\}$ would ``represent'' $\ell$ as an infinite-length vector.  However, correctly defining ``basis'' for infinite-dimensional spaces is rather difficult!  The best kind of ``basis'' is the \emph{complete orthonormal sequence} (Chapter 4), an but we are not there yet; also see the footnote on page 11.  Problem \textbf{P1} only requires that you figure out how to do some integrals, and you don't have to prove anything is a basis.  We will eventually see that if we change the vector space to a Hilbert space then this $\ell$ \emph{does} have a representation, and it is represented either by the function $e^x$ itself or the coefficients on a complete orthonormal sequence.

\prob{P2}  Show that $a_k = 1/k$, $k\in\NN$, is in $\ell^2$, and find its norm.  Is it in $\ell^1$?  Is it in $\ell^\infty$?

% P3 on A2?:  apply trapezoid rule to do the integrals in part (c), to get $\tilde \ell_k$.  for m=20 subintervals, compute $\|\ell
\end{document}
